\documentclass[10pt]{article}

% ============================================================
% LaTeX for Research Papers — Extended Quick Writing Reference
% Focus: commands used to WRITE an academic research paper.
% No field-specific econometrics / macro / theory formulas.
% ============================================================

\usepackage[landscape,margin=0.42in,includefoot]{geometry}
\usepackage{multicol}
\usepackage{amsmath,amssymb,amsthm,mathtools,bm}
\usepackage{booktabs,threeparttable,tabularx,array,multirow}
\usepackage{graphicx,subcaption}
\usepackage{enumitem}
\usepackage{xcolor}
\usepackage[hidelinks]{hyperref}
\usepackage{microtype}
\usepackage{listings}
\usepackage{float}
\usepackage{fancyhdr}
\usepackage{setspace}
\usepackage{pdflscape}
\usepackage{placeins}
\usepackage{adjustbox}
\usepackage{siunitx}
\usepackage{url}

% ---------- compact page ----------
\setlength{\parindent}{0pt}
\setlength{\parskip}{1.8pt}
\setlength{\columnsep}{0.30in}
\setlength{\columnseprule}{0.15pt}
\setlength{\tabcolsep}{4pt}
\renewcommand{\arraystretch}{1.10}
\setlist{nosep,leftmargin=*}
\raggedcolumns
\emergencystretch=1.5em
\urlstyle{same}

% ---------- code ----------
\lstset{
  basicstyle=\ttfamily\scriptsize,
  columns=fullflexible,
  breaklines=true,
  breakatwhitespace=false,
  keepspaces=true,
  frame=single,
  framerule=0.2pt,
  framesep=2.5pt,
  aboveskip=2.5pt,
  belowskip=2.5pt,
  showstringspaces=false,
  xleftmargin=1pt,
  xrightmargin=1pt
}

% ---------- useful shortcuts used in examples ----------
\newcommand{\E}{\mathbb{E}}
\newcommand{\R}{\mathbb{R}}
\newcommand{\dd}{\,\mathrm{d}}

% ---------- visual helpers ----------
\definecolor{rulegray}{gray}{0.72}
\newcommand{\topichead}[1]{%
  \vspace{3pt}\noindent\textbf{\large #1}\par\vspace{1pt}%
  {\color{rulegray}\hrule}\vspace{2pt}}
\newcommand{\shead}[1]{\vspace{1pt}\textbf{#1}\par}
\newcommand{\code}[1]{\texttt{\detokenize{#1}}}
\newcommand{\good}[1]{\textbf{#1}}

\pagestyle{fancy}
\fancyhf{}
\lfoot{LaTeX for Research Papers --- Extended Quick Writing Reference}
\rfoot{\thepage}
\renewcommand{\headrulewidth}{0pt}
\renewcommand{\footrulewidth}{0pt}

\begin{document}

\begin{center}
{\LARGE\bfseries LaTeX for Research Papers --- Extended Quick Writing Reference}\\[-1pt]
{\small Paper-writing commands, reusable source blocks, and compiled examples --- designed for quick lookup}
\end{center}
\vspace{-3pt}
\footnotesize

\begin{multicols}{2}

% ============================================================
\topichead{1. Minimal paper setup}
% ============================================================

\begin{lstlisting}
\documentclass[11pt]{article}
\usepackage[margin=1in]{geometry}
\usepackage{amsmath,amssymb,mathtools,bm}
\usepackage{booktabs,threeparttable}
\usepackage{graphicx,subcaption}
\usepackage{natbib}
\usepackage[hidelinks]{hyperref}
\usepackage{microtype}
\usepackage{setspace}

\begin{document}
...
\end{document}
\end{lstlisting}

\begin{tabularx}{\linewidth}{@{}lX@{}}
\toprule
Package & What it does in a paper \\
\midrule
\code{geometry} & margins / page size \\
\code{amsmath} & equations, \code{align}, \code{cases} \\
\code{amssymb} & extra mathematical symbols \\
\code{mathtools} & extensions to \code{amsmath} \\
\code{bm} & bold mathematical symbols \\
\code{booktabs} & professional horizontal table rules \\
\code{threeparttable} & table notes below tables \\
\code{graphicx} & figures and image scaling \\
\code{subcaption} & multi-panel figures \\
\code{natbib} & author--year citations \\
\code{hyperref} & clickable links and references \\
\code{microtype} & improves line breaking / typography \\
\code{setspace} & one-half or double spacing \\
\bottomrule
\end{tabularx}

% ============================================================
\topichead{2. A practical paper preamble}
% ============================================================

\begin{lstlisting}
\documentclass[11pt]{article}

% Page and typography
\usepackage[margin=1in]{geometry}
\usepackage{microtype}
\usepackage{setspace}

% Math
\usepackage{amsmath,amssymb,mathtools,bm}

% Tables and figures
\usepackage{booktabs,threeparttable}
\usepackage{graphicx,subcaption}

% References and links
\usepackage{natbib}
\usepackage[hidelinks]{hyperref}

% Optional conveniences
\usepackage{enumitem}
\usepackage{xcolor}
\end{lstlisting}

\textbf{Good rule:} begin with a small preamble. Add packages only when you actually need them.

% ============================================================
\topichead{3. Title, author, affiliation, date}
% ============================================================

\begin{lstlisting}
\title{Paper Title}
\author{
  First Author\thanks{University A. Email: ...}
  \and
  Second Author\thanks{University B. Email: ...}
}
\date{September 2026}

\begin{document}
\maketitle
\end{lstlisting}

Common date choices:
\begin{lstlisting}
\date{\today}          % automatic
\date{}                % suppress date
\date{September 2026}  % fixed
\end{lstlisting}

\textbf{Anonymous draft:}
\begin{lstlisting}
\author{Anonymous}
\date{}
\end{lstlisting}

% ============================================================
\topichead{4. Abstract, keywords, and JEL lines}
% ============================================================

\begin{lstlisting}
\begin{abstract}
This paper studies ...
\end{abstract}

\medskip
\noindent\textbf{Keywords:} keyword one;
keyword two; keyword three.

\noindent\textbf{JEL Codes:} A10, C10, D10.
\end{lstlisting}

Compiled style:

\textbf{Keywords:} keyword one; keyword two; keyword three.

\textbf{JEL Codes:} A10, C10, D10.

% ============================================================
\topichead{5. Normal paper section structure}
% ============================================================

\begin{lstlisting}
\section{Introduction}
\section{Background}
\section{Data}
\section{Method}
\section{Results}
\section{Conclusion}

\subsection{Sample Construction}
\subsubsection{Additional Details}
\end{lstlisting}

Starred versions suppress numbering:
\begin{lstlisting}
\section*{Acknowledgments}
\subsection*{Data Availability}
\end{lstlisting}

Compiled hierarchy:

\textbf{1 Introduction}

\hspace{1em}\textbf{1.1 Motivation}

\hspace{2em}\textbf{1.1.1 Additional detail}

% ============================================================
\topichead{6. Paragraphs, line breaks, and spacing}
% ============================================================

\begin{tabularx}{\linewidth}{@{}lX@{}}
\toprule
Source & Use \\
\midrule
blank line & new paragraph \\
\code{\\} & forced line break; use sparingly \\
\code{\noindent} & no paragraph indent \\
\code{\smallskip} & small vertical gap \\
\code{\medskip} & medium vertical gap \\
\code{\bigskip} & large vertical gap \\
\code{\newpage} & immediately start new page \\
\code{\clearpage} & print pending floats, then new page \\
\bottomrule
\end{tabularx}

Draft spacing:
\begin{lstlisting}
\onehalfspacing
\doublespacing
\singlespacing
\end{lstlisting}

% ============================================================
\topichead{7. Text formatting you actually use}
% ============================================================

\begin{tabularx}{\linewidth}{@{}Xl@{}}
\toprule
Source & Compiled \\
\midrule
\code{\textbf{important}} & \textbf{important} \\
\code{\textit{term}} & \textit{term} \\
\code{\emph{emphasis}} & \emph{emphasis} \\
\code{\texttt{file.tex}} & \texttt{file.tex} \\
\code{\underline{text}} & \underline{text} \\
\code{\textsc{small caps}} & \textsc{small caps} \\
\bottomrule
\end{tabularx}

Usually prefer \code{\emph{}} over underlining.

% ============================================================
\topichead{8. Quotes, hyphens, dashes, and special characters}
% ============================================================

\begin{lstlisting}
``quoted text''
short-run              % hyphen
1990--2000              % en dash
A result---with a note  % em dash
\end{lstlisting}

Special characters:
\begin{center}
\code{\%} $\to$ \% \quad
\code{\$} $\to$ \$ \quad
\code{\&} $\to$ \& \quad
\code{\_} $\to$ \_ \quad
\code{\#} $\to$ \#
\end{center}

Nonbreaking space:
\begin{lstlisting}
Table~\ref{tab:main}
Figure~\ref{fig:main}
Section~\ref{sec:data}
Author~(2026)
\end{lstlisting}

The \code{~} prevents an ugly line break between the word and its number.

% ============================================================
\topichead{9. Footnotes and block quotations}
% ============================================================

Footnote:
\begin{lstlisting}
The sample is restricted to adults.\footnote{
Additional restrictions are described in Appendix A.}
\end{lstlisting}

Block quotation:
\begin{lstlisting}
\begin{quote}
A short quotation or indented passage.
\end{quote}
\end{lstlisting}

Compiled:
\begin{quote}
A short quotation or indented passage.
\end{quote}

% ============================================================
\topichead{10. Labels and cross-references}
% ============================================================

\begin{lstlisting}
\section{Data}\label{sec:data}

\begin{equation}
  a=b+c.
  \label{eq:identity}
\end{equation}

\begin{table}[htbp]
  \caption{Summary Statistics}
  \label{tab:summary}
  ...
\end{table}

Section~\ref{sec:data}
Equation~\eqref{eq:identity}
Table~\ref{tab:summary}
\end{lstlisting}

Consistent prefixes save time:
\begin{center}
\code{sec:} \quad \code{eq:} \quad \code{tab:} \quad \code{fig:} \quad \code{app:}
\end{center}

% ============================================================
\topichead{11. Inline math vs. displayed math}
% ============================================================

\begin{lstlisting}
Inline math: $x_i+y_i=z_i$ appears in a sentence.

Displayed math:
\[
x_i+y_i=z_i.
\]
\end{lstlisting}

Compiled:

Inline math: \(x_i+y_i=z_i\) appears in a sentence.

\[
x_i+y_i=z_i.
\]

Use displayed math only when the expression deserves visual emphasis.

% ============================================================
\topichead{12. High-frequency math syntax}
% ============================================================

\begin{tabularx}{\linewidth}{@{}XlXl@{}}
\toprule
Source & Result & Source & Result \\
\midrule
\code{x_i} & \(x_i\) & \code{x_{it}} & \(x_{it}\) \\
\code{x^2} & \(x^2\) & \code{x^{t+1}} & \(x^{t+1}\) \\
\code{\hat{\theta}} & \(\hat\theta\) & \code{\bar{x}} & \(\bar x\) \\
\code{\tilde{x}} & \(\tilde x\) & \code{\dot{x}} & \(\dot x\) \\
\code{\frac{a}{b}} & \(\frac ab\) & \code{\sqrt{x}} & \(\sqrt x\) \\
\code{\sum_{i=1}^n x_i} & \(\sum_{i=1}^n x_i\) & \code{\prod_{i=1}^n x_i} & \(\prod_{i=1}^n x_i\) \\
\code{\int_a^b f(x)\,dx} & \(\int_a^b f(x)\,dx\) & \code{\lim_{n\to\infty}} & \(\lim_{n\to\infty}\) \\
\bottomrule
\end{tabularx}

% ============================================================
\topichead{13. Greek letters and common symbols}
% ============================================================

\begin{tabularx}{\linewidth}{@{}XlXlXl@{}}
\toprule
Source & Result & Source & Result & Source & Result \\
\midrule
\code{\alpha} & \(\alpha\) & \code{\beta} & \(\beta\) & \code{\gamma} & \(\gamma\) \\
\code{\delta} & \(\delta\) & \code{\theta} & \(\theta\) & \code{\lambda} & \(\lambda\) \\
\code{\mu} & \(\mu\) & \code{\sigma} & \(\sigma\) & \code{\rho} & \(\rho\) \\
\code{\phi} & \(\phi\) & \code{\epsilon} & \(\epsilon\) & \code{\omega} & \(\omega\) \\
\code{\le} & \(\le\) & \code{\ge} & \(\ge\) & \code{\neq} & \(\neq\) \\
\code{\approx} & \(\approx\) & \code{\in} & \(\in\) & \code{\notin} & \(\notin\) \\
\code{\to} & \(\to\) & \code{\Rightarrow} & \(\Rightarrow\) & \code{\iff} & \(\iff\) \\
\bottomrule
\end{tabularx}

% ============================================================
\topichead{14. Fonts and operators inside equations}
% ============================================================

\begin{tabularx}{\linewidth}{@{}Xl@{}}
\toprule
Source & Compiled \\
\midrule
\code{\mathbf{x}} & \(\mathbf{x}\) \\
\code{\bm{\theta}} & \(\bm{\theta}\) \\
\code{\mathbb{R}} & \(\mathbb{R}\) \\
\code{\mathcal{F}} & \(\mathcal{F}\) \\
\code{\mathrm{Var}(X)} & \(\mathrm{Var}(X)\) \\
\code{\operatorname{Cov}(X,Y)} & \(\operatorname{Cov}(X,Y)\) \\
\code{\text{if } x>0} & \(\text{if }x>0\) \\
\bottomrule
\end{tabularx}

Use \code{\text{...}} for normal prose inside math mode.

% ============================================================
\topichead{15. Align: the default for multi-line equations}
% ============================================================

\begin{lstlisting}
\begin{align}
A &= B+C, \label{eq:first}\\
  &= D+E, \\
  &= F.   \label{eq:last}
\end{align}
\end{lstlisting}

Compiled:
\begin{align}
A &= B+C, \label{eq:first-demo}\\
  &= D+E, \\
  &= F.   \label{eq:last-demo}
\end{align}

Use \code{align*} when you do not want equation numbers.

% ============================================================
\topichead{16. Long equations: split and aligned}
% ============================================================

One equation number:
\begin{lstlisting}
\begin{equation}
\begin{split}
A ={}& B+C+D+E \\
    & +F+G+H.
\end{split}
\end{equation}
\end{lstlisting}

Several aligned lines inside display math:
\begin{lstlisting}
\[
\begin{aligned}
a &= b+c,\\
d &= e+f.
\end{aligned}
\]
\end{lstlisting}

% ============================================================
\topichead{17. Delimiters, cases, matrices, vectors}
% ============================================================

Automatic delimiter size:
\begin{lstlisting}
\left( \frac{a}{b} \right)
\left[ x+y \right]
\left\{ x:x>0 \right\}
\end{lstlisting}

Cases:
\begin{lstlisting}
f(x)=
\begin{cases}
a, & \text{if } x>0,\\
b, & \text{otherwise}.
\end{cases}
\end{lstlisting}

Matrix:
\begin{lstlisting}
\begin{pmatrix}
a & b\\
c & d
\end{pmatrix}
\end{lstlisting}

Compiled:
\[
\begin{pmatrix}a&b\\c&d\end{pmatrix},\qquad
f(x)=\begin{cases}a,&x>0,\\b,&\text{otherwise.}\end{cases}
\]

% ============================================================
\topichead{18. Math spacing and punctuation}
% ============================================================

Useful spacing commands:
\begin{center}
\code{\,} small \quad
\code{\;} medium \quad
\code{\quad} large \quad
\code{\!} negative thin space
\end{center}

Example:
\begin{lstlisting}
\int f(x)\,\mathrm{d}x
\end{lstlisting}
produces
\[
\int f(x)\,\mathrm{d}x.
\]

Equations are part of sentences, so punctuation still matters.

% ============================================================
\topichead{19. Basic paper table}
% ============================================================

\begin{lstlisting}
\begin{table}[htbp]
\centering
\caption{Summary Statistics}
\label{tab:summary}
\begin{tabular}{lrrr}
\toprule
Variable & Mean & SD & N \\
\midrule
Variable A & 1.24 & 0.53 & 1,000 \\
Variable B & 3.10 & 1.22 & 1,000 \\
\bottomrule
\end{tabular}
\end{table}
\end{lstlisting}

Compiled:
\begin{center}
\begin{tabular}{lrrr}
\toprule
Variable & Mean & SD & N \\
\midrule
Variable A & 1.24 & 0.53 & 1,000 \\
Variable B & 3.10 & 1.22 & 1,000 \\
\bottomrule
\end{tabular}
\end{center}

\textbf{Paper default:} prefer \code{booktabs}; avoid vertical rules.

% ============================================================
\topichead{20. Column alignment and widths}
% ============================================================

Column letters:
\begin{center}
\code{l} = left \quad \code{c} = center \quad \code{r} = right
\end{center}

Fixed-width paragraph column:
\begin{lstlisting}
\begin{tabular}{lp{2.6in}r}
...
\end{tabular}
\end{lstlisting}

Flexible-width table:
\begin{lstlisting}
\begin{tabularx}{\textwidth}{lXr}
...
\end{tabularx}
\end{lstlisting}

Suppress outer padding:
\begin{lstlisting}
\begin{tabular}{@{}lrr@{}}
...
\end{tabular}
\end{lstlisting}

% ============================================================
\topichead{21. Multi-column headers and table panels}
% ============================================================

\begin{lstlisting}
\begin{tabular}{lcccc}
\toprule
& \multicolumn{2}{c}{Group 1}
& \multicolumn{2}{c}{Group 2} \\
\cmidrule(lr){2-3}\cmidrule(lr){4-5}
Variable & A & B & C & D \\
\midrule
... \\
\bottomrule
\end{tabular}
\end{lstlisting}

Compiled header:
\begin{center}
\begin{tabular}{lcccc}
\toprule
& \multicolumn{2}{c}{Group 1} & \multicolumn{2}{c}{Group 2} \\
\cmidrule(lr){2-3}\cmidrule(lr){4-5}
Variable & A & B & C & D \\
\midrule
Example & 1 & 2 & 3 & 4 \\
\bottomrule
\end{tabular}
\end{center}

Panel heading:
\begin{lstlisting}
\multicolumn{5}{l}{\textit{Panel A: Main Sample}} \\
\end{lstlisting}

% ============================================================
\topichead{22. Table notes: threeparttable}
% ============================================================

\begin{lstlisting}
\begin{table}[htbp]
\centering
\caption{Main Results}
\begin{threeparttable}
\begin{tabular}{lcc}
\toprule
& (1) & (2) \\
\midrule
... \\
\bottomrule
\end{tabular}
\begin{tablenotes}[flushleft]\footnotesize
\item \textit{Notes:} This table reports ...
\item Standard errors are shown in parentheses.
\end{tablenotes}
\end{threeparttable}
\end{table}
\end{lstlisting}

This is cleaner than placing a long paragraph manually beneath the table.

% ============================================================
\topichead{23. Decimal-aligned numerical columns}
% ============================================================

Optional package:
\begin{lstlisting}
\usepackage{siunitx}
\end{lstlisting}

Then:
\begin{lstlisting}
\begin{tabular}{l S[table-format=2.2]}
\toprule
Variable & {Estimate} \\
\midrule
A & 1.25 \\
B & -0.42 \\
C & 12.30 \\
\bottomrule
\end{tabular}
\end{lstlisting}

\code{S} aligns numbers around the decimal point automatically.

% ============================================================
\topichead{24. Wide tables: scale only when needed}
% ============================================================

First try smaller text:
\begin{lstlisting}
\begin{table}[htbp]
\centering
\small
...
\end{table}
\end{lstlisting}

If necessary:
\begin{lstlisting}
\resizebox{\textwidth}{!}{%
  \begin{tabular}{lcccccc}
  ...
  \end{tabular}
}
\end{lstlisting}

Or safer with \code{adjustbox}:
\begin{lstlisting}
\begin{adjustbox}{max width=\textwidth}
\begin{tabular}{lcccccc}
...
\end{tabular}
\end{adjustbox}
\end{lstlisting}

\textbf{Prefer simplifying the table before shrinking it.}

% ============================================================
\topichead{25. Landscape table page}
% ============================================================

Preamble:
\begin{lstlisting}
\usepackage{pdflscape}
\end{lstlisting}

Body:
\begin{lstlisting}
\begin{landscape}
\begin{table}[p]
\centering
\caption{Large Table}
...
\end{table}
\end{landscape}
\end{lstlisting}

The PDF viewer rotates that page automatically.

% ============================================================
\topichead{26. Import a table generated by R / Stata / Python}
% ============================================================

If software exports a LaTeX table to \texttt{tables/main.tex}:
\begin{lstlisting}
\begin{table}[htbp]
\centering
\caption{Main Results}
\label{tab:main}
\input{tables/main.tex}
\end{table}
\end{lstlisting}

Keeping generated tables in separate files makes the manuscript much easier to maintain.

% ============================================================
\topichead{27. Basic figure}
% ============================================================

\begin{lstlisting}
\begin{figure}[htbp]
\centering
\includegraphics[width=0.80\textwidth]
  {figures/main_figure.pdf}
\caption{Main Figure}
\label{fig:main}
\end{figure}
\end{lstlisting}

Compiled layout idea:
\begin{center}
\fbox{\rule{0pt}{0.75in}\rule{0.72\linewidth}{0pt}}\\[-1pt]
\textit{Figure 1: Main Figure}
\end{center}

For line charts and plots, vector PDF is usually preferable to screenshots.

% ============================================================
\topichead{28. Figure size, crop, trim, and rotation}
% ============================================================

Common sizing:
\begin{lstlisting}
\includegraphics[width=\textwidth]{fig.pdf}
\includegraphics[width=0.75\linewidth]{fig.pdf}
\end{lstlisting}

Crop margins:
\begin{lstlisting}
\includegraphics[
  width=0.8\textwidth,
  trim={1cm 0.5cm 1cm 0.5cm},
  clip
]{fig.pdf}
\end{lstlisting}

Rotate:
\begin{lstlisting}
\includegraphics[angle=90,width=0.8\textheight]{fig.pdf}
\end{lstlisting}

Trim order is \textbf{left bottom right top}.

% ============================================================
\topichead{29. Two-panel figure}
% ============================================================

\begin{lstlisting}
\begin{figure}[htbp]
\centering
\begin{subfigure}{0.48\textwidth}
  \centering
  \includegraphics[width=\linewidth]{a.pdf}
  \caption{Panel A}
  \label{fig:a}
\end{subfigure}
\hfill
\begin{subfigure}{0.48\textwidth}
  \centering
  \includegraphics[width=\linewidth]{b.pdf}
  \caption{Panel B}
  \label{fig:b}
\end{subfigure}
\caption{Main Figure Title}
\label{fig:panels}
\end{figure}
\end{lstlisting}

Inside each panel, use \code{width=\linewidth}.

% ============================================================
\topichead{30. Four-panel figure}
% ============================================================

\begin{lstlisting}
\begin{figure}[htbp]
\centering
\begin{subfigure}{0.48\textwidth}
  \includegraphics[width=\linewidth]{a.pdf}
  \caption{A}
\end{subfigure}
\hfill
\begin{subfigure}{0.48\textwidth}
  \includegraphics[width=\linewidth]{b.pdf}
  \caption{B}
\end{subfigure}

\medskip
\begin{subfigure}{0.48\textwidth}
  \includegraphics[width=\linewidth]{c.pdf}
  \caption{C}
\end{subfigure}
\hfill
\begin{subfigure}{0.48\textwidth}
  \includegraphics[width=\linewidth]{d.pdf}
  \caption{D}
\end{subfigure}
\caption{Four-Panel Figure}
\end{figure}
\end{lstlisting}

% ============================================================
\topichead{31. Figure notes / source notes}
% ============================================================

Simple method:
\begin{lstlisting}
\begin{figure}[htbp]
\centering
\includegraphics[width=.8\textwidth]{fig.pdf}
\caption{Title}
\label{fig:title}
\begin{minipage}{0.8\textwidth}
\footnotesize
\textit{Notes:} The figure shows ...
\end{minipage}
\end{figure}
\end{lstlisting}

Use the same width for the note and the figure for a clean appearance.

% ============================================================
\topichead{32. Float placement without fighting LaTeX}
% ============================================================

\begin{tabularx}{\linewidth}{@{}lX@{}}
\toprule
Option & Meaning \\
\midrule
\code{h} & approximately here \\
\code{t} & top of page \\
\code{b} & bottom of page \\
\code{p} & dedicated float page \\
\code{!} & relax placement restrictions \\
\code{H} & exactly here; requires \code{float} \\
\bottomrule
\end{tabularx}

Good defaults:
\begin{lstlisting}
\begin{table}[htbp]
\begin{figure}[htbp]
\end{lstlisting}

To prevent floats from crossing a boundary:
\begin{lstlisting}
\usepackage{placeins}
...
\FloatBarrier
\end{lstlisting}

% ============================================================
\topichead{33. Citations with natbib}
% ============================================================

\begin{tabularx}{\linewidth}{@{}Xl@{}}
\toprule
Source & Typical output \\
\midrule
\code{\citet{smith2026}} & Smith (2026) \\
\code{\citep{smith2026}} & (Smith, 2026) \\
\code{\citep[p.~5]{smith2026}} & (Smith, 2026, p. 5) \\
\code{\citeauthor{smith2026}} & Smith \\
\code{\citeyear{smith2026}} & 2026 \\
\code{\citealp{smith2026}} & Smith, 2026 \\
\bottomrule
\end{tabularx}

Multiple citations:
\begin{lstlisting}
\citep{smith2026,jones2024,lee2023}
\end{lstlisting}

Text before / after citation:
\begin{lstlisting}
\citep[see][p.~12]{smith2026}
\end{lstlisting}

% ============================================================
\topichead{34. Citation patterns used in prose}
% ============================================================

\begin{lstlisting}
\citet{smith2026} shows that ...

Related work reaches similar conclusions
\citep{jones2024,lee2023}.

This mechanism is discussed by
\citet[p.~18]{smith2026}.
\end{lstlisting}

Compiled style:

Smith (2026) shows that ...

Related work reaches similar conclusions (Jones, 2024; Lee, 2023).

% ============================================================
\topichead{35. Bibliography commands}
% ============================================================

At the end of the paper:
\begin{lstlisting}
\bibliographystyle{plainnat}
\bibliography{references}
\end{lstlisting}

Your project then contains:
\begin{lstlisting}
main.tex
references.bib
\end{lstlisting}

Many journals provide their own bibliography style. If so, replace \code{plainnat} with the required style rather than manually formatting entries.

% ============================================================
\topichead{36. BibTeX: journal article}
% ============================================================

\begin{lstlisting}
@article{smith2026,
  author  = {Smith, Jane and Lee, John},
  title   = {Article Title},
  journal = {Journal Name},
  year    = {2026},
  volume  = {10},
  number  = {2},
  pages   = {100--130},
  doi     = {10.xxxx/xxxxx}
}
\end{lstlisting}

Keep capitalization that BibTeX must preserve inside braces:
\begin{lstlisting}
title = {The Effects of {AI} on Productivity}
\end{lstlisting}

% ============================================================
\topichead{37. BibTeX: book, chapter, working paper}
% ============================================================

Book:
\begin{lstlisting}
@book{bookkey,
  author    = {Last, First},
  title     = {Book Title},
  year      = {2026},
  publisher = {Publisher}
}
\end{lstlisting}

Chapter:
\begin{lstlisting}
@incollection{chapterkey,
  author    = {Last, First},
  title     = {Chapter Title},
  booktitle = {Book Title},
  editor    = {Editor, Name},
  year      = {2026},
  publisher = {Publisher},
  pages     = {1--25}
}
\end{lstlisting}

Working paper:
\begin{lstlisting}
@techreport{wpkey,
  author      = {Last, First},
  title       = {Working Paper Title},
  institution = {NBER},
  year        = {2026},
  number      = {12345}
}
\end{lstlisting}

% ============================================================
\topichead{38. URLs, DOI links, and project links}
% ============================================================

\begin{lstlisting}
\url{https://example.com}

\href{https://example.com}{Project website}

\href{https://doi.org/10.xxxx/xxxx}{DOI}
\end{lstlisting}

If a URL causes an overfull line, use \code{\url{...}} rather than typing the URL as ordinary text.

% ============================================================
\topichead{39. Hyperref: PDF metadata}
% ============================================================

Optional:
\begin{lstlisting}
\hypersetup{
  pdftitle  = {Paper Title},
  pdfauthor = {Your Name},
  pdfsubject = {Economics},
  pdfkeywords = {keyword1, keyword2}
}
\end{lstlisting}

This affects PDF metadata, not the visible title page.

% ============================================================
\topichead{40. Acknowledgments and data statements}
% ============================================================

\begin{lstlisting}
\section*{Acknowledgments}
I thank ...

\section*{Data Availability}
Replication files are available at ...

\section*{Conflict of Interest}
The authors declare ...
\end{lstlisting}

Use only the statements required by your journal, department, or data provider.

% ============================================================
\topichead{41. Lists inside a paper}
% ============================================================

Bullets:
\begin{lstlisting}
\begin{itemize}
\item First contribution.
\item Second contribution.
\end{itemize}
\end{lstlisting}

Numbered list:
\begin{lstlisting}
\begin{enumerate}
\item First step.
\item Second step.
\end{enumerate}
\end{lstlisting}

Custom labels:
\begin{lstlisting}
\begin{enumerate}[label=(\alph*)]
\item ...
\item ...
\end{enumerate}
\end{lstlisting}

% ============================================================
\topichead{42. Theorem / proposition / lemma / assumption}
% ============================================================

Preamble:
\begin{lstlisting}
\newtheorem{theorem}{Theorem}
\newtheorem{proposition}{Proposition}
\newtheorem{lemma}{Lemma}
\newtheorem{assumption}{Assumption}
\end{lstlisting}

Body:
\begin{lstlisting}
\begin{proposition}
State the proposition here.
\end{proposition}

\begin{proof}
Write the proof here.
\end{proof}
\end{lstlisting}

Useful in theory papers and technical appendices.

% ============================================================
\topichead{43. Definitions and remarks with amsthm}
% ============================================================

\begin{lstlisting}
\theoremstyle{definition}
\newtheorem{definition}{Definition}

\theoremstyle{remark}
\newtheorem*{remark}{Remark}
\end{lstlisting}

Then:
\begin{lstlisting}
\begin{definition}
A definition goes here.
\end{definition}

\begin{remark}
An unnumbered remark goes here.
\end{remark}
\end{lstlisting}

% ============================================================
\topichead{44. Appendix basics}
% ============================================================

\begin{lstlisting}
\appendix

\section{Additional Results}
\label{app:results}

\section{Proofs}
\label{app:proofs}
\end{lstlisting}

Then reference it as:
\begin{lstlisting}
See Appendix~\ref{app:results}.
\end{lstlisting}

Place \code{\appendix} after the bibliography if that matches your preferred manuscript order.

% ============================================================
\topichead{45. Appendix equations, tables, and figures}
% ============================================================

After \code{\appendix}, standard section numbering usually becomes A, B, C, etc. For explicit appendix-style float numbering, you can reset counters:
\begin{lstlisting}
\setcounter{table}{0}
\renewcommand{\thetable}{A\arabic{table}}

\setcounter{figure}{0}
\renewcommand{\thefigure}{A\arabic{figure}}
\end{lstlisting}

For multiple appendix sections, more elaborate numbering may be handled by the journal template.

% ============================================================
\topichead{46. Large projects: separate files}
% ============================================================

Suggested structure:
\begin{lstlisting}
paper/
  main.tex
  references.bib
  sections/
    intro.tex
    background.tex
    data.tex
    results.tex
  tables/
    summary.tex
  figures/
    figure1.pdf
\end{lstlisting}

Inside \code{main.tex}:
\begin{lstlisting}
\input{sections/intro}
\input{sections/background}
\input{sections/data}
\input{sections/results}
\end{lstlisting}

No \code{\begin{document}} is needed inside the imported section files.

% ============================================================
\topichead{47. input vs. include}
% ============================================================

\begin{tabularx}{\linewidth}{@{}lX@{}}
\toprule
Command & Practical use \\
\midrule
\code{\input{file}} & insert content exactly here; best default for paper sections/tables \\
\code{\include{file}} & starts on a new page and creates auxiliary files; useful for book-like projects \\
\bottomrule
\end{tabularx}

For most research manuscripts, \code{\input} is simpler.

% ============================================================
\topichead{48. Reusable custom commands}
% ============================================================

\begin{lstlisting}
\newcommand{\E}{\mathbb{E}}
\newcommand{\R}{\mathbb{R}}
\newcommand{\vect}[1]{\bm{#1}}
\newcommand{\secref}[1]{Section~\ref{#1}}
\newcommand{\figref}[1]{Figure~\ref{#1}}
\newcommand{\tabref}[1]{Table~\ref{#1}}
\end{lstlisting}

Then:
\begin{lstlisting}
$\E[X]$
$X\in\R$
\vect{x}
\figref{fig:main}
\end{lstlisting}

Use custom commands for notation or wording you expect to change globally.

% ============================================================
\topichead{49. Draft comments and visible TODOs}
% ============================================================

Invisible source comment:
\begin{lstlisting}
% Rewrite this paragraph.
% Check citation.
\end{lstlisting}

Visible draft note:
\begin{lstlisting}
\textcolor{red}{TODO: Add citation.}
\end{lstlisting}

Reusable command:
\begin{lstlisting}
\newcommand{\todo}[1]{%
  \textcolor{red}{[TODO: #1]}}
\end{lstlisting}

Before submission, search the source for \texttt{TODO}, \texttt{FIXME}, question marks, and colored notes.

% ============================================================
\topichead{50. Temporarily hide a large block}
% ============================================================

Without another package:
\begin{lstlisting}
\iffalse
This entire block is ignored by LaTeX.
It can contain multiple paragraphs.
\fi
\end{lstlisting}

Useful when testing whether a section, figure, or table is causing an error.

% ============================================================
\topichead{51. Page layout and paragraph style}
% ============================================================

Margins:
\begin{lstlisting}
\usepackage[margin=1in]{geometry}
\end{lstlisting}

Traditional indented paragraphs:
\begin{lstlisting}
\setlength{\parindent}{1.5em}
\setlength{\parskip}{0pt}
\end{lstlisting}

Block-style paragraphs:
\begin{lstlisting}
\setlength{\parindent}{0pt}
\setlength{\parskip}{6pt}
\end{lstlisting}

Do not combine large paragraph indentation with large paragraph spacing unless a template explicitly requests it.

% ============================================================
\topichead{52. Page numbering and title-page control}
% ============================================================

Normal page numbers:
\begin{lstlisting}
\pagestyle{plain}
\end{lstlisting}

Hide only the current page number:
\begin{lstlisting}
\thispagestyle{empty}
\end{lstlisting}

Restart page numbering:
\begin{lstlisting}
\setcounter{page}{1}
\pagenumbering{arabic}
\end{lstlisting}

Roman numerals, when needed:
\begin{lstlisting}
\pagenumbering{roman}
\end{lstlisting}

% ============================================================
\topichead{53. Submission drafts: double spacing}
% ============================================================

\begin{lstlisting}
\usepackage{setspace}
...
\doublespacing
\end{lstlisting}

Or only a portion:
\begin{lstlisting}
\begin{spacing}{1.5}
Text here.
\end{spacing}
\end{lstlisting}

Use the journal or instructor specification if one exists.

% ============================================================
\topichead{54. Submission drafts: line numbers}
% ============================================================

Optional package:
\begin{lstlisting}
\usepackage{lineno}
\end{lstlisting}

Body:
\begin{lstlisting}
\linenumbers
\end{lstlisting}

Turn them off later with:
\begin{lstlisting}
\nolinenumbers
\end{lstlisting}

Some journal submission systems request line-numbered manuscripts.

% ============================================================
\topichead{55. Anonymous vs. named versions}
% ============================================================

A simple manual switch:
\begin{lstlisting}
% Anonymous submission
\author{Anonymous}

% Working-paper version
% \author{Your Name\\University}
\end{lstlisting}

For a long-running project, a Boolean switch or separate title file can reduce mistakes between versions.

% ============================================================
\topichead{56. Keep a figure or table near discussion}
% ============================================================

Usually let LaTeX float it. If you need a boundary:
\begin{lstlisting}
\usepackage{placeins}
...
End of results discussion.
\FloatBarrier
\section{Robustness}
\end{lstlisting}

If you instead use \code{[H]} everywhere, page layout often becomes worse.

% ============================================================
\topichead{57. Text sizes inside tables and figures}
% ============================================================

From larger to smaller:
\begin{center}
\code{\normalsize}\quad
\code{\small}\quad
\code{\footnotesize}\quad
\code{\scriptsize}
\end{center}

Example:
\begin{lstlisting}
\begin{table}[htbp]
\centering
\small
...
\end{table}
\end{lstlisting}

Avoid \code{\tiny} in a main paper unless there is no reasonable alternative.

% ============================================================
\topichead{58. Centering: environment vs. command}
% ============================================================

For floats:
\begin{lstlisting}
\begin{figure}[htbp]
\centering
...
\end{figure}
\end{lstlisting}

For standalone material:
\begin{lstlisting}
\begin{center}
Centered text.
\end{center}
\end{lstlisting}

Inside a float, \code{\centering} is generally cleaner than nesting a \code{center} environment.

% ============================================================
\topichead{59. Minipage for side-by-side text or notes}
% ============================================================

\begin{lstlisting}
\begin{minipage}[t]{0.48\textwidth}
Left block of text.
\end{minipage}
\hfill
\begin{minipage}[t]{0.48\textwidth}
Right block of text.
\end{minipage}
\end{lstlisting}

Useful for custom title pages, figure notes, or side-by-side nonfloating material.

% ============================================================
\topichead{60. Prevent awkward line breaks}
% ============================================================

Use nonbreaking spaces:
\begin{lstlisting}
Figure~1
Table~2
Section~3
p.~10
10~percent
\end{lstlisting}

If a particular phrase should stay together:
\begin{lstlisting}
\mbox{New York}
\end{lstlisting}

Use \code{\mbox} sparingly because it disables line breaking inside the box.

% ============================================================
\topichead{61. Common table / figure caption pattern}
% ============================================================

\begin{lstlisting}
\caption{Short Informative Title}
\label{tab:descriptive-name}
\end{lstlisting}

Put \code{\label} immediately after \code{\caption} for tables and figures. This avoids cross-reference errors.

% ============================================================
\topichead{62. A clean table-note convention}
% ============================================================

A common pattern is:
\begin{lstlisting}
\textit{Notes:} This table reports ...
Sample restrictions are ...
Standard errors are ...
Significance markers are ...
\end{lstlisting}

The note should explain the table sufficiently without forcing the reader to search the main text for basic definitions.

% ============================================================
\topichead{63. A clean figure-note convention}
% ============================================================

\begin{lstlisting}
\textit{Notes:} The figure plots ...
The shaded region indicates ...
The sample includes ...
Source: ...
\end{lstlisting}

Keep implementation detail out of the caption when it belongs in the text or appendix.

% ============================================================
\topichead{64. Referencing several items in one sentence}
% ============================================================

\begin{lstlisting}
Tables~\ref{tab:a} and~\ref{tab:b}

Figures~\ref{fig:a}--\ref{fig:c}

Sections~\ref{sec:data} and~\ref{sec:results}
\end{lstlisting}

The double hyphen \code{--} renders an en dash for a range.

% ============================================================
\topichead{65. Equation references inside prose}
% ============================================================

\begin{lstlisting}
Equation~\eqref{eq:main} defines the object.

Using \eqref{eq:first} and \eqref{eq:second},
we obtain ...
\end{lstlisting}

Use \code{\eqref} rather than manually typing parentheses around \code{\ref}.

% ============================================================
\topichead{66. Text inside equations}
% ============================================================

\begin{lstlisting}
\[
f(x)=0 \quad \text{for all } x\in\mathcal{X}.
\]
\end{lstlisting}

Compiled:
\[
f(x)=0 \quad \text{for all } x\in\mathcal{X}.
\]

Do not write ordinary words as bare math letters.

% ============================================================
\topichead{67. Underbrace / overbrace for explanation}
% ============================================================

\begin{lstlisting}
\[
A = \underbrace{B+C}_{\text{first component}}
  + \underbrace{D}_{\text{second component}}.
\]
\end{lstlisting}

Compiled:
\[
A = \underbrace{B+C}_{\text{first component}}
  + \underbrace{D}_{\text{second component}}.
\]

Useful in exposition, but avoid over-annotating final journal equations.

% ============================================================
\topichead{68. Reserved characters --- memorize these}
% ============================================================

These characters have a special meaning in LaTeX source. To print them literally, use the form in the second column.

\begin{tabularx}{\linewidth}{@{}c l X@{}}
\toprule
Character & Type this & Why it is special \\
\midrule
\# & \code{\#} & macro parameter character \\
\$ & \code{\$} & starts / ends inline math \\
\% & \code{\%} & starts a comment; rest of source line is ignored \\
\& & \code{\&} & column / alignment separator \\
\_ & \code{\_} & subscript in math mode \\
\{ & \code{\{} & begins a group / command argument \\
\} & \code{\}} & ends a group / command argument \\
\textasciitilde & \code{\textasciitilde} & \code{~} itself means a nonbreaking space \\
\textasciicircum & \code{\textasciicircum} & \code{^} is superscript syntax in math \\
\textbackslash & \code{\textbackslash} & begins a LaTeX command \\
\bottomrule
\end{tabularx}

\textbf{Most common paper examples:}
\begin{lstlisting}
10\% of the sample
R\&D expenditure
The fee is \$500.
data\_clean.csv
Figure \#2
\{A,B,C\}
\end{lstlisting}

\textbf{Dangerous one:} an unescaped percent sign comments out everything to the right of it in the \emph{source}, often causing mysterious missing text or a later brace error.

% ============================================================
\topichead{69. Literal code, filenames, URLs, and special text}
% ============================================================

For a short literal string, \code{\verb} is often easier than escaping every character:
\begin{lstlisting}
\verb|data_final_v2.csv|
\verb|beta_hat = x_1/x_2|
\end{lstlisting}

For URLs and paths:
\begin{lstlisting}
\url{https://example.com/a_b?x=1&y=2}
\path{tables/main_results_v3.tex}
\end{lstlisting}

These commands handle many special characters for you. Do \textbf{not} use \code{\verb} inside section titles, captions, footnotes, or other command arguments.

For literal angle brackets in prose, use:
\begin{lstlisting}
\textless  value \textgreater
\end{lstlisting}

% ============================================================
\topichead{70. Quotes, dashes, ellipses, and spaces}
% ============================================================

\begin{tabularx}{\linewidth}{@{}l l X@{}}
\toprule
Want & Type & Compiled idea \\
\midrule
English quotes & \code{``quoted text''} & ``quoted text'' \\
Hyphen & \code{high-frequency} & high-frequency \\
Number range & \code{1990--2000} & 1990--2000 \\
Em dash & \code{result---however} & result---however \\
Ellipsis & \code{\ldots} & \ldots \\
No line break & \code{Table~\ref{tab:main}} & keep label with number \\
Forced normal space & \code{Dr.\ Smith} & prevents sentence-sized spacing \\
\bottomrule
\end{tabularx}

\textbf{Useful habits in papers:}
\begin{lstlisting}
Figure~\ref{fig:event}
Table~\ref{tab:main}
Section~\ref{sec:data}
Appendix~\ref{app:proofs}
Dr.\ Smith
U.S.\ data
\end{lstlisting}

Plain LaTeX source using \code{`` ''}, \code{--}, and \code{---} is usually more portable than pasted typographic quote or dash characters.

% ============================================================
\topichead{71. Common writing mistakes in LaTeX papers}
% ============================================================

\begin{tabularx}{\linewidth}{@{}X X@{}}
\toprule
Avoid & Prefer \\
\midrule
\code{\\} after every paragraph & blank line between paragraphs \\
repeated line breaks for vertical space & sectioning or controlled spacing \\
\code{$$ ... $$} & \code{\[ ... \]} or \code{equation} \\
\code{$for all x$} & \code{\text{for all } x} inside math \\
manual ``Figure 3'' & \code{Figure~\ref{fig:key}} \\
manual ``Table 2'' & \code{Table~\ref{tab:key}} \\
label before figure/table caption & \code{\caption} then \code{\label} \\
shrinking a whole table immediately & fix columns / headings first \\
using \code{[H]} for every float & start with \code{[htbp]} \\
spaces to align text & table, \code{align}, or proper environments \\
manual bibliography entries for a large paper & BibTeX / BibLaTeX database \\
\bottomrule
\end{tabularx}

\textbf{Paragraph rule:} a double backslash means ``line break,'' not ``new paragraph.'' A blank source line starts a new paragraph.

% ============================================================
\topichead{72. Common equation-writing mistakes}
% ============================================================

An equation is grammatically part of the sentence. Punctuate it accordingly.

\textbf{Better:}
\begin{lstlisting}
The objective can be written as
\begin{equation}
Q(\theta)=\sum_{i=1}^{n} q_i(\theta),
\label{eq:objective}
\end{equation}
where $\theta$ denotes the parameter vector.
\end{lstlisting}

\textbf{Common mistakes:}
\begin{itemize}
\item Use \code{\text{...}} for ordinary words inside math.
\item Use built-in operators such as \code{\log}, \code{\exp}, \code{\sin}, and \code{\max}; do not type them as ordinary italic letters.
\item Do not mix unmatched \code{$}, \code{\(}, \code{\)}, \code{\[}, or \code{\]} delimiters.
\item Avoid putting long prose inside display math just to center it.
\item Put \code{\label} inside the numbered equation environment you want to reference.
\end{itemize}

\topichead{73. Common compile sequence with BibTeX}
% ============================================================

Classic workflow:
\begin{lstlisting}
pdflatex main.tex
bibtex main
pdflatex main.tex
pdflatex main.tex
\end{lstlisting}

Why multiple passes? LaTeX needs later passes to resolve citations, equation numbers, figure numbers, table numbers, and page references.

% ============================================================
\topichead{74. Easier compiling with latexmk}
% ============================================================

If available:
\begin{lstlisting}
latexmk -pdf main.tex
\end{lstlisting}

Clean auxiliary files:
\begin{lstlisting}
latexmk -c
\end{lstlisting}

Full cleanup including PDF:
\begin{lstlisting}
latexmk -C
\end{lstlisting}

Editors such as Overleaf and VS Code + LaTeX Workshop usually automate the repeated passes.

% ============================================================
\topichead{75. Fast debugging: common errors}
% ============================================================

\begin{tabularx}{\linewidth}{@{}lX@{}}
\toprule
Symptom & Check first \\
\midrule
Undefined control sequence & typo, wrong command name, or missing package \\
Missing \$ inserted & math command / underscore / caret outside math mode \\
Extra alignment tab \& & unescaped \code{&} or too many table/alignment columns \\
Missing \} inserted & unmatched brace, often earlier than the reported line \\
Runaway argument? & missing closing brace or environment end \\
Environment mismatch & wrong / missing \code{\end{...}} \\
Reference shows \code{??} & compile again; verify unique label and \code{\ref} key \\
Citation shows \code{?} & wrong citation key, missing \texttt{.bib}, or BibTeX pass \\
Multiply-defined label & same \code{\label} used more than once \\
File \code{...sty} not found & package not installed / wrong package name \\
File \code{...} not found & image/input path, capitalization, or extension \\
Overfull \code{\hbox} & long URL, equation, table, or unbreakable text \\
Underfull \code{\hbox} & awkward line stretching; often only a warning \\
Float(s) lost / misplaced & float nested illegally or placement too restrictive \\
\bottomrule
\end{tabularx}

\textbf{Best debugging rule: fix the \emph{first} reported LaTeX error first.} Later errors are often consequences of that first one.

\textbf{If an error seems unrelated to the reported line:} inspect several lines \emph{above} it for an unescaped percent sign, unmatched \code{$}, missing brace, or missing \code{\end{...}}.

% ============================================================
\topichead{76. Overfull lines: safe fixes}
% ============================================================

For URLs:
\begin{lstlisting}
\url{https://very-long-address.example/...}
\end{lstlisting}

For equations: use \code{align}, \code{split}, or \code{multline} rather than shrinking the font.

For tables: shorten headings, use \code{tabularx}, reduce columns, or move a very wide table to landscape.

Last-resort local relaxation:
\begin{lstlisting}
\begin{sloppypar}
A difficult paragraph ...
\end{sloppypar}
\end{lstlisting}

% ============================================================
\topichead{77. Figure / table appears too late}
% ============================================================

Try, in this order:
\begin{enumerate}
\item use \code{[htbp]};
\item reduce the float's size;
\item move its source slightly earlier;
\item use \code{\FloatBarrier} at a logical boundary;
\item use \code{\clearpage} before an appendix or references if appropriate.
\end{enumerate}

Avoid solving every float problem with \code{[H]}.

% ============================================================
\topichead{78. Full copy-ready paper skeleton}
% ============================================================

\begin{lstlisting}
\documentclass[11pt]{article}
\usepackage[margin=1in]{geometry}
\usepackage{amsmath,amssymb,mathtools,bm}
\usepackage{booktabs,threeparttable}
\usepackage{graphicx,subcaption}
\usepackage{natbib}
\usepackage[hidelinks]{hyperref}
\usepackage{microtype}
\usepackage{setspace}

\title{Paper Title}
\author{Your Name\\University Name}
\date{\today}

\begin{document}
\maketitle

\begin{abstract}
One-paragraph abstract.
\end{abstract}

\noindent\textbf{Keywords:} ...

\noindent\textbf{JEL Codes:} ...

\section{Introduction}\label{sec:intro}
\input{sections/intro}

\section{Background}\label{sec:background}
\input{sections/background}

\section{Data}\label{sec:data}
\input{sections/data}

\section{Method}\label{sec:method}
\input{sections/method}

\section{Results}\label{sec:results}
\input{sections/results}

\section{Conclusion}\label{sec:conclusion}
\input{sections/conclusion}

\bibliographystyle{plainnat}
\bibliography{references}

\appendix
\section{Additional Material}
\input{sections/appendix}

\end{document}
\end{lstlisting}

% ============================================================
\topichead{79. Useful optional packages by task}
% ============================================================

\begin{tabularx}{\linewidth}{@{}lX@{}}
\toprule
Task & Package \\
\midrule
Landscape pages & \code{pdflscape} \\
Float barriers & \code{placeins} \\
Exact float placement & \code{float} \\
Decimal alignment & \code{siunitx} \\
Fit wide object to page & \code{adjustbox} \\
Multiple table rows & \code{multirow} \\
Better lists & \code{enumitem} \\
Colored draft notes & \code{xcolor} \\
Line numbers & \code{lineno} \\
Subfigures & \code{subcaption} \\
\bottomrule
\end{tabularx}

Do not load optional packages unless you need the feature.

% ============================================================
\topichead{80. Commands to recognize instantly}
% ============================================================

\begin{lstlisting}
\documentclass   \usepackage
\begin           \end
\title           \author
\date            \maketitle
\begin{abstract} \section
\subsection      \textbf
\textit          \emph
\footnote        \label
\ref             \eqref
$...$            \[...\]
_                ^
\frac            \sqrt
\sum             \int
\hat             \bar
\mathbb          \bm
\left            \right
\text            \operatorname
\begin{align}    \begin{cases}
\begin{table}    \toprule
\midrule         \bottomrule
\multicolumn     \cmidrule
\caption         \includegraphics
\begin{figure}   \citet
\citep           \bibliography
\input           \appendix
\newcommand      \url
\href            \clearpage
\FloatBarrier    \resizebox
\end{lstlisting}

% ============================================================
\topichead{81. What to memorize vs. look up}
% ============================================================

\textbf{Memorize:} sections, text emphasis, math mode, subscripts/superscripts, \code{frac}, \code{align}, \code{label/ref}, basic tables, basic figures, \code{caption}, \code{citet/citep}, \code{input}, and \code{appendix}.

\medskip
\textbf{Look up when needed:} complex multi-panel tables, unusual symbols, custom theorem numbering, journal-specific classes, landscape material, long tables, TikZ, bibliography-style edge cases, and specialized submission formatting.

\begin{quote}
\textbf{Fast workflow: copy a working block, change one thing, compile, repeat.}
\end{quote}

% ============================================================
\topichead{82. Final paper-writing sanity check}
% ============================================================

Before sending a draft, check:
\begin{itemize}
\item every figure and table is referenced in the text;
\item captions and notes explain what the reader needs;
\item no references or citations show \code{??} or \code{?};
\item no visible TODO / FIXME notes remain;
\item figure fonts are readable at 100\% PDF zoom;
\item tables do not run outside margins;
\item bibliography entries are complete and consistent;
\item the anonymous version really removes identifying information;
\item appendix labels and references work;
\item the PDF opens cleanly after a full compile from scratch.
\end{itemize}

\end{multicols}
\end{document}
